\فصل{نتایج آزمایش‌ها}

در این فصل، نتیجهٔ مسیرهای اصلی معرفی‌شده در فصل پیش بررسی می‌شود: ریزتنظیم ویسپر-کوچک با داده‌های معمولی فارسی، ریزتنظیم همان مدل با ترکیب داده‌های معمولی و تخریب‌شده، دو اجرای فست کانفورمر، مدل همجوشی دوجریانه‌ای، ریزتنظیم ویسپر-متوسط برای فارسی، و آزمایش نهایی ویسپر-کوچک با داده‌های افزودهٔ ایران‌صدا. هدف فقط مرتب‌کردن مدل‌ها بر اساس یک عدد نهایی نیست؛ نتایج باید روشن کنند که آیا شبیه‌سازی تخریب‌های مخابراتی واقعاً به کانال‌های واقعی تعمیم می‌یابد، آیا پیچیدگی مدل همجوشی سودی فراتر از آموزش چندشرایطی دارد، افزایش ظرفیت مدل چه نتیجه‌ای در ارزیابی ثابت این پژوهش داشته است، و افزودن دادهٔ شبه‌برچسب‌خوردهٔ صوت بلند چه اثری بر دامنه‌های مختلف گفتار می‌گذارد.

بر همین اساس، تحلیل فصل حول شش پرسش پژوهش انجام می‌شود: نخست، افزودن داده‌های تخریب‌شده چه مقدار از خطای ویسپر-کوچک می‌کاهد؟ دوم، آیا این بهبود به گفتار واقعی تلفنی و \lr{VoIP} منتقل می‌شود، بی‌آن‌که در گفتار عمومی پسرفت نامقبولی ایجاد کند؟ سوم، آیا معماری بهساز-همجوشی از خط مبنای چندشرایطی بهتر است؟ چهارم، آیا دروازهٔ آموخته‌شده واقعاً از هر دو نمای اصلی و بهسازی‌شده استفاده می‌کند؟ پنجم، آیا افزودن دادهٔ شبه‌برچسب‌خوردهٔ ایران‌صدا به ترکیب معمولی، بلند و تخریب‌شده عملکرد ویسپر-کوچک را بهبود می‌دهد؟ و ششم، سود یا زیان این دادهٔ افزوده میان گفتار عمومی و گفتار واقعی تلفنی/\lr{VoIP} چگونه توزیع می‌شود؟ هر شش پرسش بر پایهٔ ارزیابی‌های تکمیل‌شده در این فصل پاسخ داده می‌شوند.

\قسمت{پروتکل ارزیابی}

دو معیار اصلی، نرخ خطای واژه\پاورقی{Word Error Rate (WER)} و نرخ خطای نویسه\پاورقی{Character Error Rate (CER)} هستند. اگر تعداد جانشینی‌ها، حذف‌ها و درج‌ها را به‌ترتیب با $S$، $D$ و $I$ و تعداد واحدهای متن مرجع را با $N$ نشان دهیم، نرخ خطا از رابطهٔ زیر به دست می‌آید:

\begin{equation}
\mathrm{ER}=\frac{S+D+I}{N}.
\end{equation}

در محاسبهٔ WER، واحد شمارش واژه و در محاسبهٔ CER، نویسه است. بنابراین مقدار کمتر در هر دو معیار بهتر است، اما این دو معیار جنبه‌های یکسانی از خطا را اندازه نمی‌گیرند. WER به درست‌بودن کل واژه حساس است، درحالی‌که CER شدت تفاوت نوشتاری درون واژه را نیز نشان می‌دهد. به همین دلیل، هم‌جهت‌نبودن این دو معیار می‌تواند نشانهٔ تفاوت در الگوی خطا باشد و نباید یکی را جایگزین کامل دیگری دانست.

همهٔ مدل‌ها با رمزگشایی حریصانه و یک زنجیرهٔ یکسان برای نرمال‌سازی متن ارزیابی شدند. پنج مجموعهٔ آزمون جدول~\رجوع{جدول:داده-آزمون} هم گفتار عمومی فارسی و هم گفتار واقعی مخابراتی را پوشش می‌دهند. AGFarsdat نرمال‌شده از دید هدف پژوهش مهم‌ترین مجموعه است، زیرا نمونه‌های تلفنی و \lr{VoIP} واقعی دارد و خروجی شبیه‌ساز تخریب این پژوهش نیست. در نتیجه، بهبود روی این مجموعه شاهدی برای انتقال آموخته‌های حاصل از تخریب مصنوعی به یک کانال واقعی است. سایر مجموعه‌ها برای بررسی این موضوع به کار می‌روند که افزایش مقاومت مخابراتی به قیمت افت گسترده در گفتار عمومی تمام نشده باشد.

سه اجرای تک‌جریانهٔ ویسپر-کوچک، شامل مدل معمولی، مدل چندشرایطی و مدل چندشرایطی با ایران‌صدا، روی $21{,}848$ نمونهٔ واجد شرایط مشترک امتیازدهی شدند. ویسپر-متوسط نیز روی همین مجموعه سنجیده شد. مدل همجوشی و دو اجرای فست کانفورمر هر $21{,}854$ نمونه را پردازش کردند؛ بنابراین، مقدارهای تجمیعی آن‌ها کاملاً با چهار مدل دیگر جفت‌شده نیست. افزایش ظرفیت ویسپر-متوسط نیز اجازه نمی‌دهد اختلاف آن با مدل‌های کوچک به‌عنوان یک مقایسهٔ صرفاً داده‌محور تفسیر شود. تفاوت‌های کوچک میان مدل همجوشی و مدل چندشرایطی نیز بدون آزمون آماری به‌عنوان برتری قطعی یا معنادار تفسیر نمی‌شوند.

تنظیمات اصلی سه اجرای تک‌جریانهٔ ویسپر-کوچک یکسان بود: بذر تصادفی $1337$، یک دورهٔ آموزش، نرخ یادگیری $10^{-5}$، اندازهٔ دستهٔ $4$ و انباشت گرادیان در $8$ گام. در اجرای چندشرایطی، دو مجموعهٔ تخریب‌شدهٔ کامن‌وویس به داده‌های معمولی و گونه‌های بلند افزوده شدند؛ اجرای ایران‌صدا همین ترکیب کامل را حفظ کرد و فقط قطعه‌های پذیرفته‌شدهٔ کتاب صوتی را به آن افزود. ازاین‌رو، مدل چندشرایطی خط مبنای مستقیم برای سنجش اثر دادهٔ ایران‌صدا است. مدل همجوشی با برنامهٔ سه‌مرحله‌ای فصل پیش آموزش یافت: حداکثر $20{,}000$ گام گرم‌کردن بهساز، $30{,}000$ گام آموزش بهساز و همجوشی با ویسپر ثابت، و $120{,}000$ گام آموزش مشترک انتهابه‌انتها. مقایسهٔ مدل همجوشی افزون بر تغییر معماری، تفاوت برنامه و مدت آموزش را نیز در بر دارد.

\قسمت{نتایج تجمیعی}

جدول~\رجوع{جدول:نتایج-کلی} نتیجهٔ تجمیعی را روی همهٔ مجموعه‌های آزمون نشان می‌دهد. ستون آخر تغییر مطلق WER نسبت به ویسپر معمولی را گزارش می‌کند؛ علامت منفی به معنای کاهش خطا است. مقدار تجمیعی میانگین سادهٔ پنج سطر جدول تفکیکی نیست و از کل پیش‌بینی‌ها و مراجع ارزیابی محاسبه شده است. در نتیجه، مجموعه‌های بزرگ‌تر و جمله‌های دارای واژه‌های بیش‌تر سهم بیشتری در این عدد دارند.

\شروع{لوح}[ht]
\تنظیم‌ازوسط
\شرح{نتایج تجمیعی روش‌های اصلی روی همهٔ مجموعه‌های آزمون}
\resizebox{\textwidth}{!}{%
\شروع{جدول}{|c|c|c|c|}
\خط‌پر
\سیاه روش & \سیاه WER & \سیاه CER & \سیاه تغییر مطلق WER نسبت به ویسپر کوچک (معمولی) \\
\خط‌پر \خط‌پر
ویسپر کوچک (معمولی) & \lr{23.27\%} & \lr{10.84\%} & --- \\
ویسپر کوچک (چندشرایطی) & \lr{17.83\%} & \lr{7.93\%} & \lr{-5.44} واحد درصد \\
همجوشی دوجریانه & \lr{18.44\%} & \lr{14.04\%} & \lr{-4.82} واحد درصد \\
ویسپر متوسط & \مهم{\lr{14.36\%}} & \مهم{\lr{6.68\%}} & \مهم{\lr{-8.91}} واحد درصد \\
ویسپر کوچک (چندشرایطی + ایران‌صدا) & \lr{17.22\%} & \lr{7.70\%} & \lr{-6.04} واحد درصد \\
\خط‌پر
\پایان{جدول}}
\برچسب{جدول:نتایج-کلی}
\پایان{لوح}

افزودن داده‌های تخریب‌شده، WER را از $23.27\%$ به $17.83\%$ رسانده است. این تغییر برابر با $5.44$ واحد درصد کاهش مطلق و حدود $23.4\%$ کاهش نسبی خطاست. CER نیز از $10.84\%$ به $7.93\%$ کاهش یافته است؛ یعنی $2.91$ واحد درصد کاهش مطلق و حدود $26.8\%$ کاهش نسبی. چون داده‌های ارزیابی و تنظیمات بهینه‌سازی دو اجرا یکسان‌اند، این نتیجه قوی‌ترین شاهد آزمایش حاضر برای اثر آموزش چندشرایطی است: مواجه‌کردن مستقیم بازشناس با نویز، محدودیت پهنای‌باند، مصنوعات کدک و افت بسته، هم در سطح واژه و هم در سطح نویسه خطا را به‌طور محسوسی کم کرده است.

مدل همجوشی نیز نسبت به ویسپر معمولی WER کمتری دارد: کاهش مطلق $4.82$ واحد درصد، معادل حدود $20.7\%$ کاهش نسبی. بااین‌حال، WER آن $0.62$ واحد درصد از مدل چندشرایطی بیشتر است. افزون بر این، CER مدل همجوشی به $14.04\%$ رسیده و از مدل معمولی و مدل چندشرایطی بدون ایران‌صدا بدتر است. بنابراین، نتیجهٔ تجمیعی دو پیام متفاوت دارد: مدل دوجریانه بخش مهمی از مقاومت مدل چندشرایطی را حفظ کرده، اما پیچیدگی بهسازی و همجوشی به بهبود کلی نسبت به خط مبنای چندشرایطی منجر نشده و الگوی خطای نویسه‌ای آن نیز نیازمند بررسی دقیق‌تر است.

افزودن ایران‌صدا به خط مبنای چندشرایطی، WER را از $17.83\%$ به $17.22\%$ کاهش داده است؛ یعنی $0.60$ واحد درصد کاهش مطلق و حدود $3.4\%$ کاهش نسبی. CER نیز از $7.93\%$ به $7.70\%$ رسیده است که $0.23$ واحد درصد کاهش مطلق و حدود $2.9\%$ کاهش نسبی است. هم‌جهت‌بودن دو معیار نشان می‌دهد دادهٔ افزوده در نتیجهٔ تجمیعی سودمند بوده است، اما اندازهٔ سود در مقایسه با اثر افزودن داده‌های تخریب‌شده محدودتر است و بدون اجرای چندبذری یا فاصلهٔ اطمینان نباید به‌عنوان برتری آماری تفسیر شود.

ویسپر-متوسط با WER برابر $14.36\%$ و CER برابر $6.68\%$ قوی‌ترین نتیجهٔ کلی را ثبت کرده است. در مقایسه با ویسپر-کوچک چندشرایطی، WER آن $3.47$ واحد درصد و CER آن $1.25$ واحد درصد کمتر است. یکسان‌بودن $21{,}848$ نمونهٔ ارزیابی، اثر اختلاف مجموعهٔ آزمون را حذف می‌کند، اما چون ظرفیت مدل نیز تغییر کرده است این فاصله را نمی‌توان فقط به ترکیب داده یا یک عامل منفرد نسبت داد. ازاین‌رو، ویسپر-متوسط بهترین مدل از نظر نتیجهٔ کلی است، درحالی‌که مقایسهٔ دو اجرای ویسپر-کوچک چندشرایطی اثر افزودن ایران‌صدا را با کنترل مناسب‌تری نشان می‌دهد.

\قسمت{نتایج به تفکیک مجموعهٔ آزمون}

عدد تجمیعی برای انتخاب مدل نهایی مفید است، اما تفاوت حوزه‌های صوتی را پنهان می‌کند. جدول \رجوع{جدول:نتایج-wer} WER را به تفکیک مجموعه نشان می‌دهد. بهترین مقدار هر سطر پررنگ شده است؛ این پررنگ‌سازی صرفاً کمترین عدد مشاهده‌شده را نشان می‌دهد و به معنای اثبات معناداری آماری اختلاف نیست.

\شروع{لوح}[ht]
\تنظیم‌ازوسط
\شرح{WER به تفکیک مجموعهٔ آزمون؛ مقدار کمتر بهتر است}
\resizebox{\textwidth}{!}{%
\شروع{جدول}{|c|c|c|c|c|c|}
\خط‌پر
\سیاه مجموعهٔ آزمون & \سیاه ویسپر کوچک (معمولی) & \سیاه ویسپر کوچک (چندشرایطی) & \سیاه همجوشی دوجریانه & \سیاه ویسپر متوسط & \سیاه ویسپر کوچک (چندشرایطی + ایران‌صدا) \\
\خط‌پر \خط‌پر
AGFarsdat (تلفنی/VoIP واقعی) & \lr{35.63\%} & \lr{25.80\%} & \lr{25.38\%} & \مهم{\lr{21.99\%}} & \lr{26.66\%} \\
کامن‌وویس فارسی ۲۵ & \lr{11.52\%} & \lr{7.84\%} & \lr{9.06\%} & \مهم{\lr{3.84\%}} & \lr{8.19\%} \\
فلورز فارسی & \lr{21.05\%} & \lr{19.82\%} & \lr{19.61\%} & \lr{17.14\%} & \مهم{\lr{13.78\%}} \\
پرشین‌اسپیچ & \lr{34.53\%} & \lr{35.25\%} & \lr{31.18\%} & \lr{31.65\%} & \مهم{\lr{26.86\%}} \\
پیکرهٔ گفتار فارسی & \lr{30.17\%} & \lr{30.90\%} & \lr{35.17\%} & \lr{31.17\%} & \مهم{\lr{25.87\%}} \\
\خط‌پر
\پایان{جدول}}
\برچسب{جدول:نتایج-wer}
\پایان{لوح}

\زیرقسمت{گفتار واقعی تلفنی و VoIP}

مهم‌ترین بهبود در AGFarsdat دیده می‌شود. آموزش چندشرایطی WER را از $35.63\%$ به $25.80\%$ کاهش داده است؛ یعنی $9.83$ واحد درصد کاهش مطلق و حدود $27.6\%$ کاهش نسبی. چون AGFarsdat در تولید دادهٔ تخریب‌شده یا آموزش مدل نقشی نداشته و کانال‌های واقعی مخابراتی را شامل می‌شود، این نتیجه نشان می‌دهد تنوع مصنوعی آموزش صرفاً حفظ نشده، بلکه دست‌کم تا حدی به شرایط واقعی انتقال یافته است.

مدل همجوشی در همین مجموعه WER برابر $25.38\%$ دارد و در میان مدل‌های کوچکِ بدون ایران‌صدا کمترین مقدار را به دست آورده است. اختلاف آن با مدل چندشرایطی فقط $0.42$ واحد درصد، یا حدود $1.6\%$ نسبت به خطای مدل چندشرایطی، است. بدون بازنمونه‌گیری آماری یا فاصلهٔ اطمینان نمی‌توان این فاصلهٔ کوچک را برتری پایدار مدل همجوشی دانست. تفسیر محتاطانه‌تر آن است که مدل دوجریانه در دامنهٔ اصلی پژوهش عملکرد مدل چندشرایطی را حفظ کرده و اندکی بهبود عددی ایجاد کرده است، اما این بهبود کوچک برای جبران افت آن در برخی مجموعه‌های بزرگ‌تر کافی نیست.

مدل دارای ایران‌صدا در AGFarsdat به WER برابر $26.66\%$ رسیده است؛ یعنی نسبت به خط مبنای چندشرایطی $0.86$ واحد درصد پسرفت دارد. بنابراین، بهبود تجمیعی این مدل به گفتار واقعی تلفنی و \lr{VoIP} منتقل نشده است. بااین‌حال، خطای آن همچنان $8.97$ واحد درصد کمتر از ویسپر معمولی است؛ در نتیجه، افزودن ایران‌صدا بخش عمدهٔ مقاومت حاصل از آموزش چندشرایطی را حفظ کرده، اما آن را در دامنهٔ مخابراتی تقویت نکرده است.

\زیرقسمت{گفتار عمومی و تفاوت حوزه‌ها}

در کامن‌وویس، مدل چندشرایطی WER را از $11.52\%$ به $7.84\%$ می‌رساند؛ کاهش مطلق $3.68$ واحد درصد و کاهش نسبی حدود $31.9\%$. بخشی از این سود می‌تواند ناشی از آن باشد که دادهٔ تخریب‌شده از بخش آموزش همین پیکره ساخته شده است، هرچند نمونه‌های آزمون در آموزش حضور نداشته‌اند. مدل همجوشی با WER برابر $9.06\%$ همچنان از ویسپر معمولی بهتر است، اما $1.23$ واحد درصد از مدل چندشرایطی عقب می‌ماند. افزودن ایران‌صدا نیز WER را به $8.19\%$ می‌رساند که $0.35$ واحد درصد از خط مبنای چندشرایطی بیشتر است؛ بنابراین، دادهٔ کتاب صوتی در این مجموعه سودی ایجاد نکرده است.

در فلورز، آموزش چندشرایطی WER را $1.23$ واحد درصد کاهش می‌دهد و مدل همجوشی با $19.61\%$ فقط $0.20$ واحد درصد بهتر از آن است. در مقابل، مدل دارای ایران‌صدا WER برابر $13.78\%$ ثبت می‌کند: $6.04$ واحد درصد کمتر از خط مبنای چندشرایطی و $3.36$ واحد درصد کمتر از ویسپر-متوسط. این بزرگ‌ترین بهبود قابل اتکای ایران‌صدا در یک مجموعهٔ عمومی با اندازهٔ متوسط است و نشان می‌دهد تنوع گفتار بلند شبه‌برچسب‌خورده به این دامنه منتقل شده است.

پرشین‌اسپیچ تنها $24$ نمونهٔ ارزیابی‌شده دارد. مدل دارای ایران‌صدا WER را نسبت به مدل چندشرایطی $8.39$ واحد درصد کاهش می‌دهد و با $26.86\%$ بهترین مقدار این سطر را ثبت می‌کند. بااین‌حال، هر نمونه می‌تواند اثر بزرگی بر نتیجه داشته باشد؛ بنابراین این سطر بیشتر یک مشاهدهٔ مقدماتی است تا شاهدی محکم برای تعمیم‌پذیری. برای نتیجه‌گیری معتبر، به مجموعه‌ای بزرگ‌تر از همان دامنه یا فاصلهٔ اطمینان مبتنی بر بازنمونه‌گیری نیاز است.

در پیکرهٔ گفتار فارسی، آموزش چندشرایطی به‌تنهایی WER مدل معمولی را اندکی از $30.17\%$ به $30.90\%$ افزایش می‌دهد و مدل همجوشی با $35.17\%$ ضعیف‌ترین نتیجه را ثبت می‌کند. افزودن ایران‌صدا این روند را معکوس می‌کند و WER را به $25.87\%$ می‌رساند؛ یعنی $5.03$ واحد درصد بهبود نسبت به مدل چندشرایطی و بهترین مقدار این سطر. این تغییر نشان می‌دهد گفتار بلند افزوده‌شده می‌تواند بخشی از ناهمخوانی حوزه‌ای این پیکره را جبران کند، هرچند داده‌های موجود برای نسبت‌دادن سود به ویژگی مشخصی مانند سبک گفتار، کیفیت ضبط یا توزیع متن کافی نیستند.

در مقایسه با ویسپر-کوچک چندشرایطی، ویسپر-متوسط WER را در چهار مجموعهٔ AGFarsdat، کامن‌وویس، فلورز و پرشین‌اسپیچ کاهش می‌دهد. تنها استثنا پیکرهٔ گفتار فارسی است که WER از $30.90\%$ به $31.17\%$، یعنی $0.27$ واحد درصد، افزایش یافته است. بزرگ‌ترین سود عددی در AGFarsdat دیده می‌شود: WER مدل متوسط $3.81$ واحد درصد کمتر از خط مبنای چندشرایطی است. این الگو با نتیجهٔ تجمیعی سازگار است، ولی به دلیل تغییر ظرفیت مدل شاهدی برای اثر صرف داده نیست.

\قسمت{تحلیل خطای نویسه}

برای آشکارشدن تفاوت الگوی خطا، جدول~\رجوع{جدول:نتایج-cer} مقدار CER هر مجموعه را نیز گزارش می‌کند. ویسپر-متوسط در نتیجهٔ تجمیعی بهترین CER را دارد، اما بهترین مدل در همهٔ مجموعه‌ها یکسان نیست.

\شروع{لوح}[ht]
\تنظیم‌ازوسط
\شرح{CER به تفکیک مجموعهٔ آزمون؛ مقدار کمتر بهتر است}
\resizebox{\textwidth}{!}{%
\شروع{جدول}{|c|c|c|c|c|c|}
\خط‌پر
\سیاه مجموعهٔ آزمون & \سیاه ویسپر کوچک (معمولی) & \سیاه ویسپر کوچک (چندشرایطی) & \سیاه همجوشی دوجریانه & \سیاه ویسپر متوسط & \سیاه ویسپر کوچک (چندشرایطی + ایران‌صدا) \\
\خط‌پر \خط‌پر
AGFarsdat (تلفنی/VoIP واقعی) & \lr{19.73\%} & \lr{12.62\%} & \lr{12.38\%} & \مهم{\lr{10.92\%}} & \lr{12.96\%} \\
کامن‌وویس فارسی ۲۵ & \lr{4.06\%} & \lr{3.22\%} & \lr{16.08\%} & \مهم{\lr{1.65\%}} & \lr{3.29\%} \\
فلورز فارسی & \lr{6.13\%} & \lr{6.36\%} & \lr{6.23\%} & \lr{5.42\%} & \مهم{\lr{4.09\%}} \\
پرشین‌اسپیچ & \lr{16.21\%} & \lr{15.62\%} & \lr{12.91\%} & \lr{13.99\%} & \مهم{\lr{9.82\%}} \\
پیکرهٔ گفتار فارسی & \lr{15.48\%} & \lr{15.22\%} & \lr{21.07\%} & \lr{17.29\%} & \مهم{\lr{14.16\%}} \\
\خط‌پر
\پایان{جدول}}
\برچسب{جدول:نتایج-cer}
\پایان{لوح}

در AGFarsdat، نتیجهٔ CER نتیجهٔ WER را تأیید می‌کند: مدل دارای ایران‌صدا با CER برابر $12.96\%$ نسبت به خط مبنای چندشرایطی $0.35$ واحد درصد پسرفت دارد. در کامن‌وویس نیز CER آن با افزایش $0.08$ واحد درصدی به $3.29\%$ می‌رسد. در مقابل، ایران‌صدا CER را در فلورز از $6.36\%$ به $4.09\%$، در پرشین‌اسپیچ از $15.62\%$ به $9.82\%$ و در پیکرهٔ گفتار فارسی از $15.22\%$ به $14.16\%$ کاهش می‌دهد. هم‌جهت‌بودن تغییر WER و CER در هر پنج مجموعه نشان می‌دهد سود و زیان ایران‌صدا به حوزه وابسته است، نه این‌که فقط از تفاوت واحد اندازه‌گیری ناشی شده باشد.

CER مدل همجوشی روی کامن‌وویس به $16.08\%$ می‌رسد، با آن‌که WER همان اجرا $9.06\%$ است. همچنین CER آن در پیکرهٔ گفتار فارسی $21.07\%$ است. همین دو مجموعه بخش مهمی از افزایش CER تجمیعی مدل همجوشی را توضیح می‌دهند.

از فایل‌های معیار تجمیعی نمی‌توان تعیین کرد که این واگرایی دقیقاً از چند خروجی بسیار بلند، تکرار نویسه‌ها، خطاهای نرمال‌سازی یا الگوی دیگری ناشی شده است. نتیجهٔ قابل دفاع در وضعیت فعلی این است که WER نسبتاً خوب مدل همجوشی نباید ضعف CER آن را پنهان کند. تحلیل نمونه‌به‌نمونهٔ درج، حذف و جانشینی و بازبینی خروجی‌های پرت، پیش‌نیاز هر ادعای دقیق‌تر دربارهٔ علت این رفتار است.

\قسمت{نتایج خط مبنای فست کانفورمر}

\شروع{لوح}[ht]
\تنظیم‌ازوسط
\شرح{نتایج فست کانفورمر با داده‌های معمولی و چندشرایطی}
\resizebox{\textwidth}{!}{%
\شروع{جدول}{|c|c|c|c|c|}
\خط‌پر
\سیاه مجموعهٔ آزمون & \سیاه WER معمولی & \سیاه CER معمولی & \سیاه WER چندشرایطی & \سیاه CER چندشرایطی \\
\خط‌پر \خط‌پر
تجمیعی & \lr{22.73\%} & \lr{14.79\%} & \lr{21.81\%} & \lr{14.45\%} \\
AGFarsdat (تلفنی/VoIP واقعی) & \lr{31.07\%} & \lr{14.81\%} & \lr{30.48\%} & \lr{14.09\%} \\
کامن‌وویس فارسی ۲۵ & \lr{11.68\%} & \lr{16.46\%} & \lr{10.01\%} & \lr{16.10\%} \\
فلورز فارسی & \lr{29.22\%} & \lr{9.94\%} & \lr{29.40\%} & \lr{10.22\%} \\
پرشین‌اسپیچ & \lr{30.70\%} & \lr{13.50\%} & \lr{31.41\%} & \lr{13.56\%} \\
پیکرهٔ گفتار فارسی & \lr{33.69\%} & \lr{12.22\%} & \lr{33.59\%} & \lr{12.68\%} \\
\خط‌پر
\پایان{جدول}}
\برچسب{جدول:نتایج-فست-کانفورمر}
\پایان{لوح}

آموزش چندشرایطی در فست کانفورمر نیز نتیجهٔ کلی را بهبود داد. WER تجمیعی از $22.73\%$ به $21.81\%$ کاهش یافت که برابر با $0.92$ واحد درصد کاهش مطلق و حدود $4.1\%$ کاهش نسبی است. CER نیز از $14.79\%$ به $14.45\%$ رسید. در AGFarsdat، WER از $31.07\%$ به $30.48\%$ و CER از $14.81\%$ به $14.09\%$ کاهش یافت.

بیشترین بهبود WER این مدل در کامن‌وویس دیده شد، درحالی‌که در فلورز و پرشین‌اسپیچ تغییر کوچکی در جهت نامطلوب رخ داد. بنابراین، افزودن داده‌های تخریب‌شده برای فست کانفورمر نیز در نتیجهٔ تجمیعی و در گفتار واقعی مخابراتی سودمند بود، اما اندازهٔ بهبود آن از ویسپر-کوچک کمتر و توزیع آن میان مجموعه‌ها ناهمگون‌تر بود.

مقایسهٔ مستقیم اعداد فست کانفورمر و ویسپر باید با احتیاط انجام شود، زیرا این دو مدل از نظر پیش‌آموزش، معماری رمزگشا و تنظیمات آموزش یکسان نیستند. در اینجا نتیجهٔ اصلی از مقایسهٔ دو اجرای فست کانفورمر با یکدیگر به دست می‌آید.

\قسمت{تهدیدهای اعتبار و دامنهٔ نتیجه‌گیری}

نتایج این فصل شواهد روشنی به سود آموزش چندشرایطی فراهم می‌کنند، اما دامنهٔ ادعا باید با محدودیت‌های آزمایش سازگار بماند:

\شروع{فقرات}
\فقره هر پیکربندی تک‌جریانه با یک بذر و یک اجرای آموزشی گزارش شده است. در نتیجه، پراکندگی ناشی از مقداردهی اولیه، ترتیب دسته‌ها و انتخاب نقطهٔ وارسی اندازه‌گیری نشده است.
\فقره برای اختلاف مدل‌ها فاصلهٔ اطمینان یا آزمون معناداری آماری محاسبه نشده است. این محدودیت به‌ویژه برای برتری‌های کوچک همجوشی در AGFarsdat و فلورز و بهبود تجمیعی محدود مدل ایران‌صدا مهم است.
\فقره مدل همجوشی شش نمونهٔ بیش‌تر از سه مدل ویسپر-کوچک ارزیابی کرده است. بنابراین مقایسهٔ تجمیعی آن کاملاً جفت‌شده نیست، هرچند اندازهٔ اختلاف نمونه کوچک است.
\فقره ویسپر-متوسط روی همان نمونه‌های سه ویسپر-کوچک ارزیابی شده است، اما ظرفیت بیش‌تر آن مانع از تفسیر اختلاف به‌عنوان اثر کنترل‌شدهٔ ترکیب داده می‌شود.
\فقره ویسپر و فست کانفورمر از نظر معماری، پیش‌آموزش، رمزگشایی و تنظیمات بهینه‌سازی متفاوت‌اند. بنابراین، اختلاف مطلق آن‌ها را نمی‌توان تنها به نوع معماری نسبت داد و مقایسهٔ اصلی برای فست کانفورمر میان دو اجرای معمولی و چندشرایطی همان مدل انجام می‌شود.
\فقره WER و CER به تفکیک نمایهٔ تخریب، بازهٔ SNR، کدک، نرخ افت بسته و طول کلیپ گزارش نشده‌اند. در نتیجه، معلوم نیست کدام جزء شبیه‌ساز بیشترین سهم را در مقاومت ایجاد کرده است.
\فقره نتیجهٔ تجمیعی تحت تأثیر مجموعه‌های بزرگ AGFarsdat و کامن‌وویس است و نباید به معنای عملکرد یکنواخت در همهٔ دامنه‌ها تفسیر شود.
\فقره اثر ایران‌صدا میان مجموعه‌ها ناهمگون است: سه مجموعهٔ عمومی بهبود یافته‌اند، اما AGFarsdat و کامن‌وویس پسرفت کوچکی دارند. بدون تحلیل ویژگی‌های صوتی و متنی قطعه‌های افزوده‌شده، علت این تفاوت حوزه‌ای مشخص نیست.
\فقره ارزیابی حاضر کیفیت بازشناسی را اندازه می‌گیرد، نه کیفیت ادراکی صوت یا بازنمایی بهسازی‌شده. معیارهایی مانند SI-SDR و PESQ و نیز آزمون شنیداری اجرا نشده‌اند.
\فقره هزینهٔ محاسباتی، حافظهٔ مصرفی و تأخیر استنتاج مدل دوجریانه اندازه‌گیری نشده است. ازاین‌رو، حتی در مجموعه‌هایی که همجوشی اندکی بهتر است، دربارهٔ صرفهٔ استقرار آن نمی‌توان نتیجه‌گیری کرد.
\پایان{فقرات}

\قسمت{جمع‌بندی فصل}

پاسخ پرسش نخست روشن است: افزودن داده‌های تخریب‌شده به آموزش ویسپر-کوچک، WER تجمیعی را $23.4\%$ و CER را $26.8\%$ به‌صورت نسبی کاهش داده است. پاسخ پرسش دوم نیز مثبت اما محدود است: کاهش نسبی WER در AGFarsdat واقعی حدود $27.6\%$ است و نشان می‌دهد تخریب‌های مصنوعی آموزش قابلیت انتقال به گفتار مخابراتی واقعی را دارند. در عین حال، افت کوچک در پیکرهٔ گفتار فارسی یادآور می‌شود که این سود در همهٔ دامنه‌ها یکسان نیست.

نتیجهٔ فست کانفورمر نیز با این یافته هم‌سو بود. افزودن داده‌های تخریب‌شده WER تجمیعی این مدل را از $22.73\%$ به $21.81\%$ کاهش داد و در AGFarsdat نیز بهبود کوچکی ایجاد کرد. این نتیجه نشان می‌دهد سود داده‌های مخابراتی مصنوعی به ویسپر محدود نیست، هرچند ویسپر-کوچک در آزمایش حاضر بهرهٔ بیشتری از این داده‌ها برد.

پاسخ پرسش سوم منفی است: مدل همجوشی در نتیجهٔ تجمیعی به ویسپر-کوچک چندشرایطی نمی‌رسد و CER ضعیف‌تری دارد. پاسخ پرسش چهارم نیز منفی است: میانگین وزن $0.0039$ برای نمای بهسازی‌شده نشان می‌دهد دروازه تقریباً همیشه به نمای اصلی تکیه کرده و استفادهٔ معنادار و متوازنی از دو جریان نیاموخته است. در کنار این پاسخ‌ها، ویسپر-متوسط با WER برابر $14.36\%$ و CER برابر $6.68\%$ قوی‌ترین نتیجهٔ کلی است.

پاسخ پرسش پنجم مثبت اما محدود است: افزودن دادهٔ شبه‌برچسب‌خوردهٔ ایران‌صدا به ترکیب کامل چندشرایطی، WER تجمیعی را از $17.83\%$ به $17.22\%$ و CER را از $7.93\%$ به $7.70\%$ کاهش داده است. این تغییر به‌ترتیب معادل کاهش نسبی حدود $3.4\%$ و $2.9\%$ است. بنابراین، دادهٔ افزوده در ارزیابی کلی سودمند بوده است، اما اندازهٔ این سود کوچک‌تر از اثر آموزش چندشرایطی است و با یک اجرای تک‌بذری نمی‌توان پایداری آماری آن را اثبات کرد.

پاسخ پرسش ششم نشان می‌دهد این سود یکنواخت نیست. ایران‌صدا WER را در فلورز، پرشین‌اسپیچ و پیکرهٔ گفتار فارسی به‌ترتیب $6.04$، $8.39$ و $5.03$ واحد درصد کاهش داده است، اما در کامن‌وویس $0.35$ واحد درصد و در AGFarsdat تلفنی/\lr{VoIP} واقعی $0.86$ واحد درصد پسرفت ایجاد کرده است. CER نیز همین جهت تغییر را نشان می‌دهد. در نتیجه، سود اصلی دادهٔ کتاب صوتی در سه مجموعهٔ گفتار عمومی دیده می‌شود و به دامنهٔ مخابراتی واقعی منتقل نشده است؛ هرچند مدل ایران‌صدا همچنان نسبت به ویسپر معمولی در AGFarsdat مقاوم‌تر باقی مانده است.
